\usetikzlibrary{arrows.meta}
\begin{figure}[htbp]
  \centering
  % shift={(...)} 使用坐标系单位；勿用 xshift=..cm，否则面板间距会被放大约 1/x 倍。
  \makebox[\textwidth][c]{%
  \resizebox{1.12\textwidth}{!}{%
  \begin{tikzpicture}[
      x=0.82cm,
      y=0.60cm,
      font=\scriptsize,
      panel/.style={
        draw=black!28,
        rounded corners=1.5pt,
        line width=0.35pt
      },
      lane/.style={
        draw=black!14,
        line width=0.3pt
      },
      lanelabel/.style={
        anchor=east,
        font=\fontsize{5}{6}\selectfont,
        text=black!72,
        align=right
      },
      ctrl/.style={
        draw=violet!55!black,
        fill=violet!10,
        rounded corners=1pt,
        line width=0.45pt
      },
      taskA/.style={
        draw=blue!65!black,
        fill=blue!14,
        rounded corners=0.8pt,
        line width=0.45pt
      },
      taskB/.style={
        draw=orange!80!black,
        fill=orange!18,
        rounded corners=0.8pt,
        line width=0.45pt
      },
      taskC/.style={
        draw=teal!65!black,
        fill=teal!15,
        rounded corners=0.8pt,
        line width=0.45pt
      },
      launchbox/.style={
        draw=black!70,
        fill=black!7,
        rounded corners=0.7pt,
        line width=0.4pt
      },
      boundary/.style={
        draw=black!55,
        densely dashed,
        line width=0.45pt
      },
      dispatch/.style={
        -{Latex[length=1.3mm]},
        draw=violet!70!black,
        densely dotted,
        line width=0.5pt
      },
      dep/.style={
        -{Latex[length=1.3mm]},
        draw=blue!70!black,
        line width=0.55pt
      },
      spill/.style={
        {Latex[length=1.2mm]}-{Latex[length=1.2mm]},
        draw=red!65!black,
        line width=0.55pt
      },
      wait/.style={
        draw=black!30,
        fill=black!5,
        densely dashed,
        rounded corners=0.7pt,
        line width=0.35pt
      },
      memory/.style={
        draw=black!45,
        fill=black!6,
        rounded corners=0.8pt,
        line width=0.4pt
      }
    ]

    % ======================================================================
    % (a) Separate physical kernels
    % ======================================================================
    \begin{scope}[shift={(0,0)}]
      \draw[panel] (0.05,0.10) rectangle (7.05,7.20);

      \node[font=\small\bfseries] at (3.60,6.82)
        {(a) 多 kernel};

      \foreach \yy in {5.90,5.05,4.05,3.35,2.65} {
        \draw[lane] (1.70,\yy) -- (6.85,\yy);
      }

      \node[lanelabel] at (1.58,5.90) {Host};
      \node[lanelabel] at (1.58,5.05) {Grid ctrl.};
      \node[lanelabel] at (1.58,4.05) {SM 0};
      \node[lanelabel] at (1.58,3.35) {SM 1};
      \node[lanelabel] at (1.58,2.65) {SM $n$};
      \node[lanelabel] at (1.58,1.55) {HBM};

      \draw[launchbox] (1.75,5.69) rectangle (2.22,6.11);
      \node[font=\tiny] at (1.985,5.90) {$L_A$};
      \draw[launchbox] (3.20,5.69) rectangle (3.67,6.11);
      \node[font=\tiny] at (3.435,5.90) {$L_B$};
      \draw[launchbox] (4.75,5.69) rectangle (5.22,6.11);
      \node[font=\tiny] at (4.985,5.90) {$L_C$};

      \draw[ctrl] (1.85,4.84) rectangle (3.05,5.26);
      \node[font=\tiny] at (2.45,5.05) {grid $A$};
      \draw[ctrl] (3.35,4.84) rectangle (4.60,5.26);
      \node[font=\tiny] at (3.975,5.05) {grid $B$};
      \draw[ctrl] (4.90,4.84) rectangle (6.40,5.26);
      \node[font=\tiny] at (5.65,5.05) {grid $C$};

      \foreach \yy in {4.05,3.35,2.65} {
        \draw[taskA] (1.85,\yy-0.21) rectangle (3.05,\yy+0.21);
        \node[font=\tiny] at (2.45,\yy) {$A$};
        \draw[taskB] (3.35,\yy-0.21) rectangle (4.60,\yy+0.21);
        \node[font=\tiny] at (3.975,\yy) {$B$};
        \draw[taskC] (4.90,\yy-0.21) rectangle (6.40,\yy+0.21);
        \node[font=\tiny] at (5.65,\yy) {$C$};
      }

      \draw[boundary] (3.20,2.35) -- (3.20,5.40);
      \draw[boundary] (4.75,2.35) -- (4.75,5.40);
      \node[font=\fontsize{4.5}{5}\selectfont,fill=white,inner sep=0.5pt,text=black!65]
        at (3.20,4.52) {grid 完成};

      \draw[memory] (1.70,1.33) rectangle (6.85,1.77);
      \node[font=\fontsize{5}{6}\selectfont,text=black!70] at (4.275,1.55)
        {输入 / 中间结果 / 输出};

      \draw[spill] (3.20,1.79) -- node[right,font=\tiny] {$I_A$} (3.20,2.42);
      \draw[spill] (4.75,1.79) -- node[right,font=\tiny] {$I_B$} (4.75,2.42);

      \draw[-{Latex[length=1.5mm]},black!65]
        (1.70,0.82) -- (6.80,0.82)
        node[above left,inner sep=1pt,font=\tiny] {time};
      \node[font=\fontsize{5}{6}\selectfont\bfseries,text=black!72] at (4.275,0.38)
        {3 launches $\cdot$ 2 次全局 grid 边界};
    \end{scope}

    % ======================================================================
    % (b) Conventional compile-time operator fusion
    % ======================================================================
    \begin{scope}[shift={(7.18,0)}]
      \draw[panel] (0.05,0.10) rectangle (7.05,7.20);

      \node[font=\small\bfseries] at (3.60,6.82)
        {(b) 传统算子融合};

      \foreach \yy in {5.90,5.05,4.05,3.35,2.65} {
        \draw[lane] (1.70,\yy) -- (6.85,\yy);
      }

      \node[lanelabel] at (1.58,5.90) {Host};
      \node[lanelabel] at (1.58,5.05) {Grid ctrl.};
      \node[lanelabel] at (1.58,4.05) {SM 0};
      \node[lanelabel] at (1.58,3.35) {SM 1};
      \node[lanelabel] at (1.58,2.65) {SM $n$};
      \node[lanelabel] at (1.58,1.55) {HBM};

      \draw[launchbox] (1.75,5.69) rectangle (2.65,6.11);
      \node[font=\tiny] at (2.20,5.90) {$L_{A+B}$};
      \draw[launchbox] (4.65,5.69) rectangle (5.12,6.11);
      \node[font=\tiny] at (4.885,5.90) {$L_C$};

      \draw[ctrl] (1.85,4.84) rectangle (4.45,5.26);
      \node[font=\tiny] at (3.15,5.05) {grid $F_{A+B}$};
      \draw[ctrl] (4.80,4.84) rectangle (6.40,5.26);
      \node[font=\tiny] at (5.60,5.05) {grid $C$};

      \foreach \yy in {4.05,3.35,2.65} {
        \draw[taskA] (1.85,\yy-0.21) rectangle (2.85,\yy+0.21);
        \node[font=\tiny] at (2.35,\yy) {$A$};
        \draw[taskB] (2.85,\yy-0.21) rectangle (4.45,\yy+0.21);
        \node[font=\tiny] at (3.65,\yy) {$B$};
        \draw[taskC] (4.80,\yy-0.21) rectangle (6.40,\yy+0.21);
        \node[font=\tiny] at (5.60,\yy) {$C$};
      }

      \draw[boundary] (4.625,2.35) -- (4.625,5.40);

      \draw[memory] (1.70,1.33) rectangle (6.85,1.77);
      \node[font=\fontsize{5}{6}\selectfont,text=black!70] at (4.275,1.55)
        {输入 / 融合输出 / 最终输出};

      \draw[spill] (4.625,1.79)
        -- node[right,font=\tiny] {$I_{A+B}$} (4.625,2.42);

      \draw[-{Latex[length=1.5mm]},black!65]
        (1.70,0.82) -- (6.80,0.82)
        node[above left,inner sep=1pt,font=\tiny] {time};
      \node[font=\fontsize{5}{6}\selectfont\bfseries,text=black!72] at (4.275,0.38)
        {2 launches $\cdot$ grid 内静态融合};
    \end{scope}

    % ======================================================================
    % (c) Persistent megakernel with an on-device scheduler
    % ======================================================================
    \begin{scope}[shift={(14.36,0)}]
      \draw[panel] (0.05,0.10) rectangle (7.05,7.20);

      \node[font=\small\bfseries] at (3.60,6.82)
        {(c) Persistent Megakernel};

      \foreach \yy in {5.90,5.05,4.05,3.35,2.65} {
        \draw[lane] (1.70,\yy) -- (6.85,\yy);
      }

      \node[lanelabel] at (1.58,5.90) {Host};
      \node[lanelabel] at (1.58,5.05) {Sched. SM};
      \node[lanelabel] at (1.58,4.05) {SM 0};
      \node[lanelabel] at (1.58,3.35) {SM 1};
      \node[lanelabel] at (1.58,2.65) {SM $n$};
      \node[lanelabel] at (1.58,1.55) {HBM};

      \draw[launchbox] (1.75,5.69) rectangle (2.33,6.11);
      \node[font=\tiny] at (2.04,5.90) {$L_M$};

      \draw[
        draw=teal!55!black,
        densely dashed,
        rounded corners=1.2pt,
        line width=0.55pt
      ] (1.78,2.32) rectangle (6.45,5.34);

      \node[
        font=\fontsize{5}{6}\selectfont\bfseries,
        fill=white,
        inner sep=0.5pt,
        text=teal!55!black
      ] at (5.45,5.48) {persistent grid};

      \draw[ctrl] (1.90,4.83) rectangle (6.30,5.27);
      \node[font=\fontsize{5}{6}\selectfont] at (4.10,5.05)
        {ready queue / events / counters};

      % Worker SM 0
      \draw[taskA] (1.95,3.84) rectangle (2.68,4.26);
      \node[font=\tiny] at (2.315,4.05) {$A_0$};
      \draw[taskB] (2.83,3.84) rectangle (3.58,4.26);
      \node[font=\tiny] at (3.205,4.05) {$B_0$};
      \draw[wait] (3.65,3.84) rectangle (4.85,4.26);
      \node[font=\tiny,text=black!50] at (4.25,4.05) {wait};
      \draw[taskC] (4.92,3.84) rectangle (5.72,4.26);
      \node[font=\tiny] at (5.32,4.05) {$C_0$};

      % Worker SM 1
      \draw[taskA] (2.10,3.14) rectangle (2.85,3.56);
      \node[font=\tiny] at (2.475,3.35) {$A_1$};
      \draw[taskB] (3.00,3.14) rectangle (3.75,3.56);
      \node[font=\tiny] at (3.375,3.35) {$B_1$};
      \draw[wait] (3.82,3.14) rectangle (4.98,3.56);
      \node[font=\tiny,text=black!50] at (4.40,3.35) {wait};
      \draw[taskC] (5.05,3.14) rectangle (5.85,3.56);
      \node[font=\tiny] at (5.45,3.35) {$C_1$};

      % Worker SM n
      \draw[taskA] (2.30,2.44) rectangle (3.10,2.86);
      \node[font=\tiny] at (2.70,2.65) {$A_n$};
      \draw[taskB] (3.25,2.44) rectangle (4.30,2.86);
      \node[font=\tiny] at (3.775,2.65) {$B_n$};
      \draw[wait] (4.37,2.44) rectangle (5.10,2.86);
      \node[font=\tiny,text=black!50] at (4.735,2.65) {wait};
      \draw[taskC] (5.17,2.44) rectangle (6.00,2.86);
      \node[font=\tiny] at (5.585,2.65) {$C_n$};

      \draw[dep] (2.68,4.31) -- (2.83,4.31);
      \node[
        font=\fontsize{4.5}{5}\selectfont,
        anchor=south,
        text=blue!70!black
      ] at (2.75,4.36) {tile ready};

      \draw[dispatch] (2.85,4.82) -- (2.85,4.30);
      \draw[dispatch] (5.15,4.82) -- (5.15,3.60);

      \draw[
        draw=red!65!black,
        densely dashed,
        line width=0.55pt
      ] (4.70,2.32) -- (4.70,4.70);
      \node[
        font=\fontsize{4.5}{5}\selectfont,
        anchor=south east,
        fill=white,
        inner sep=0.5pt,
        text=red!65!black
      ] at (4.62,4.36) {all ready};

      \draw[memory] (1.70,1.33) rectangle (6.85,1.77);
      \node[font=\fontsize{5}{6}\selectfont,text=black!70] at (4.275,1.55)
        {模型状态 / KV / 输出};

      \draw[-{Latex[length=1.2mm]},draw=black!55]
        (2.20,1.78) -- (2.20,2.30);
      \draw[-{Latex[length=1.2mm]},draw=black!55]
        (6.05,2.30) -- (6.05,1.78);

      \draw[-{Latex[length=1.5mm]},black!65]
        (1.70,0.82) -- (6.80,0.82)
        node[above left,inner sep=1pt,font=\tiny] {time};
      \node[font=\fontsize{5}{6}\selectfont\bfseries,text=black!72] at (4.275,0.38)
        {1 launch $\cdot$ 设备端任务依赖};
    \end{scope}
  \end{tikzpicture}%
  }}%
  \caption{同一条 $A\!\rightarrow\!B\!\rightarrow\!C$ 算子链的三种执行映射。
  (a)~多 kernel 以物理 grid completion 排序阶段，并在边界处物化中间结果；
  (b)~传统融合消去 $A$--$B$ 物理边界，但 worker 仍执行静态融合程序，且 $C$ 前仍保留物理边界；
  (c)~persistent Megakernel 仅 launch 一次，由设备端调度维护 queue/event/counter，worker 拉取 tile 任务：
  tile-ready 允许 $B_0$ 与 $A$ 尾部重叠，而 $C$ 前的 all-ready 则表现为内部 readiness gate。
  Megakernel 中间结果是否落 HBM 取决于具体实现，图中不作约定。}
  \label{fig:megakernel-principle}
\end{figure}
